% ==========================================================
% 2. JUSTIFICACIÓN (Fase 3, Tema 3.2)
% ==========================================================

\section{Justificación}

Este estudio relaciona el tiempo diario de exposición a pantallas y los hábitos
posturales autorreportados con la frecuencia de fatiga visual y de molestias
osteomusculares en estudiantes de Ingeniería de Sistemas y Computación de la UPTC. Su
intención es medir y relacionar esas variables en una población que, según los
antecedentes revisados, no ha sido estudiada de forma específica, y obtener datos
primarios del contexto local que hoy no existen.

Es necesario hacerlo ahora porque buena parte de la evidencia disponible se produjo en
condiciones de clases virtuales durante la pandemia de la COVID-19 y en otras carreras
y poblaciones [1], [3], [5], [6], mientras que no se identificó evidencia sobre las
condiciones actuales de estudio de este programa. Cada periodo académico sin datos
propios es un periodo más en el que las molestias pueden seguir viéndose como un
efecto normal de la carrera y en el que las medidas de prevención se deciden sin
información de la propia población.

El estudio sirve en tres dimensiones. En la dimensión \textbf{teórica}, aporta
evidencia sobre una población distinta a la de los estudios previos (estudiantes de
medicina, odontología y posgrado, y docentes) [1], [2], [5]--[7], y permite observar
si las asociaciones reportadas en ellos se sostienen en estudiantes de pregrado de
ingeniería. En la dimensión \textbf{práctica}, entrega a los propios estudiantes, a la
Escuela y a Bienestar Universitario información para decidir qué medidas de ergonomía
y pausas activas vale la pena promover. La dimensión \textbf{metodológica} pesa menos:
el estudio no desarrolla un instrumento nuevo, pues se apoyará en cuestionarios ya
validados para ambos grupos de síntomas, y su aporte es un diseño transversal sencillo
que otros programas o sedes podrían replicar.

Se aborda con un diseño no experimental, transversal y correlacional, mediante una
encuesta aplicada una sola vez. Es más conveniente que una intervención porque la
pregunta es si existe relación entre las variables, no cuánto mejora una medida, y
porque permite recoger información de un número amplio de estudiantes con recursos
gratuitos y en el tiempo disponible. Deja una línea base útil para estudios
posteriores, por ejemplo uno que evalúe una intervención de pausas activas.
